% ============================================================
% CHAPTER 3 - RESEARCH METHODOLOGY
% ============================================================
\chapter{Research Methodology}
\label{ch:methodology}

\section{Introduction}

This chapter sets out how the research was conducted. Section~\ref{sec:approach}
states the research approach and the overall process. Section~\ref{sec:dataset}
describes the corpus and the pairing scheme. Sections~\ref{sec:g2p} to
\ref{sec:features} describe the three stages that convert a pair of recordings
into labelled feature patches: grapheme-to-phoneme conversion, matched-pair
forced alignment and log-mel feature extraction.
Section~\ref{sec:labelling} --- the methodological core of the dissertation ---
presents the two labelling strategies that were attempted, explains why the
first failed and derives the second. Sections~\ref{sec:split} to
\ref{sec:metrics} cover the data split, the training protocol, the architecture
comparison design and the evaluation metrics. Section~\ref{sec:tools} lists the
tools and the experimental environment, and Section~\ref{sec:ethics} addresses
ethical considerations.

\section{Research Approach}
\label{sec:approach}

The work follows an \emph{experimental, design-science} approach: an artefact
(the detection model and the surrounding system) is constructed, evaluated
against measurable criteria, diagnosed when it fails, and redesigned. This is
not incidental to the presentation --- the central result of the dissertation
was obtained precisely by the diagnose-and-redesign cycle applied to the
labelling scheme.

The process comprised the seven stages shown in
Figure~\ref{fig:methodology-overview}. The dashed arrow marks the feedback path
that produced the principal contribution: qualitative inference-time evaluation
revealed that healthy control recordings were being flagged as full of errors,
which invalidated the labelling stage and forced its redesign.

\begin{figure}[H]
\centering
\begin{tikzpicture}[node distance=8mm and 10mm]
  \node[block] (lit)   {1. Literature review \& gap analysis};
  \node[block, below=of lit] (data)  {2. Corpus selection \& pairing};
  \node[block, below=of data] (align) {3. G2P + matched-pair forced alignment};
  \node[block, below=of align] (feat) {4. Log-mel feature extraction};
  \node[block, below=of feat, fill=orange!12] (label) {5. Label derivation \\ (Method 1 $\rightarrow$ Method 2)};
  \node[block, below=of label] (train) {6. Model training \& architecture comparison};
  \node[block, below=of train] (eval)  {7. Evaluation \& deployment};

  \draw[arrow] (lit)   -- (data);
  \draw[arrow] (data)  -- (align);
  \draw[arrow] (align) -- (feat);
  \draw[arrow] (feat)  -- (label);
  \draw[arrow] (label) -- (train);
  \draw[arrow] (train) -- (eval);

  \draw[dashedarrow] (eval.east) -- ++(2.8,0) |- (label.east);
  \node[font=\scriptsize, align=center, anchor=west]
        at ([xshift=3.0cm, yshift=1.6cm]eval.east)
        {control speech\\flagged as error};
\end{tikzpicture}
\caption[Research methodology overview]{Overview of the research process. The
dashed feedback path represents the empirical finding that drove the
redesign of the labelling stage from Method~1 to Method~2.}
\label{fig:methodology-overview}
\end{figure}

\section{Overview of the Proposed Pipeline}
\label{sec:pipeline-overview}

Figure~\ref{fig:methodology-overview} describes the \emph{research process};
Figure~\ref{fig:method-pipeline} describes the \emph{method} itself. It is the
complete deep-learning pipeline developed in this work, from a pair of raw
waveforms to a trained per-phoneme error classifier, and every stage in it is
specified in detail in the sections that follow. Read left to right and top to
bottom, it comprises three bands.

\textbf{The data pipeline} (top band, Sections~\ref{sec:dataset} to
\ref{sec:labelling}) turns audio into supervised training examples. Its two
distinctive steps are the alignment of \emph{both} recordings of a pair to a
single control-derived target sequence, which is what makes phoneme $i$ of the
impaired production comparable with phoneme $i$ of the healthy one, and the
control-baseline labelling rule that derives the ground truth from that
comparison. The latter is the contribution of the dissertation and is
highlighted in the figure.

\textbf{The network} (middle band, Section~\ref{sec:network-design}) consumes
one $64 \times 32$ log-mel patch per phoneme. Two convolutional blocks reduce
the patch to a $16 \times 8 \times 32$ tensor of local spectro-temporal
features; a time-distributed flatten reinterprets that as a 16-step sequence of
256-dimensional vectors; a bidirectional LSTM integrates the sequence in both
directions; and a two-unit softmax emits $P(\text{error})$. The tensor shape is
annotated beneath each stage so the reduction can be followed end to end.

\textbf{Training and evaluation} (bottom band,
Sections~\ref{sec:training-protocol} to \ref{sec:metrics}) fit the network on
the class-balanced training partition and assess it on the untouched held-out
partition. The pipeline's outputs are the two artefacts that the deployed
system consumes --- the trained network and the pair of global normalisation
constants --- together with the evaluation evidence reported in
Chapter~\ref{ch:results}.

The single most important structural property of the figure is that the
normalisation constants $\mu$ and $\sigma$ are an \emph{output} of the
pipeline, on the same footing as the network weights. Because the features are
in decibels and therefore level-dependent, a system that scales its input
differently from training is a different model, regardless of which weights it
loads.

\begin{figure}[H]
\centering
\resizebox{\textwidth}{!}{%
\begin{tikzpicture}[font=\scriptsize]
  \tikzset{
    pb/.style   = {rectangle, rounded corners=2pt, draw=black!65, fill=blue!7,
                   text width=2.95cm, align=center, minimum height=1.30cm,
                   font=\scriptsize, inner sep=3pt},
    nb/.style   = {pb, fill=orange!10},
    tb/.style   = {pb, fill=green!10, text width=4.30cm},
    hb/.style   = {pb, fill=red!12},
    band/.style = {draw=black!55, dashed, rounded corners=4pt},
    bt/.style   = {font=\small\bfseries, anchor=west},
    ar/.style   = {-{Stealth[length=2.4mm]}, thick, draw=black!70},
    nt/.style   = {font=\tiny, text=black!60, align=center},
  }

  % ===================== BAND 1 : DATA PIPELINE =====================
  \draw[band, fill=blue!2] (-0.35,0.62) rectangle (15.05,-5.60);
  \node[bt] at (-0.15,0.25)
       {1 --- DATA PIPELINE \normalfont\scriptsize(\S\ref{sec:dataset}--\S\ref{sec:labelling})};

  \node[pb] (d1) at (1.90,-1.15)  {\textbf{paired input}\\dysarthric $x_d$, control $x_c$,\\control transcript $w$};
  \node[pb] (d2) at (5.60,-1.15)  {\textbf{preprocess}\\16\,kHz mono,\\trim silence once};
  \node[pb] (d3) at (9.30,-1.15)  {\textbf{target sequence}\\espeak-ng G2P on $w$,\\strip stress marks};
  \node[pb] (d4) at (13.00,-1.15) {\textbf{forced alignment}\\wav2vec2 CTC, \emph{both}\\recordings to same target};

  \node[pb] (d5) at (1.90,-3.95)  {\textbf{feature extraction}\\$\pm20$\,ms context, 64 mel\\bands, dB, 32 frames};
  \node[pb] (d6) at (5.60,-3.95)  {\textbf{deviation}\\$D_i = \overline{|d_i - c_i|}$\\Eq.~\eqref{eq:distance}};
  \node[hb] (d7) at (9.30,-3.95)  {\textbf{control-baseline}\\\textbf{labelling (Method 2)}\\$\tau = 1.20$\,dB, Eq.~\eqref{eq:method2}};
  \node[pb] (d8) at (13.00,-3.95) {\textbf{split \& balance}\\80/20 stratified,\\upsample train only};

  \foreach \s/\t in {d1/d2, d2/d3, d3/d4, d5/d6, d6/d7, d7/d8} \draw[ar] (\s) -- (\t);
  \draw[ar] (d4.south) -- (13.00,-2.55) -- (1.90,-2.55) -- (d5.north);

  \node[nt, anchor=west] at (-0.15,-5.25)
       {98{,}306 phoneme patches $\cdot$ 14{,}746 error / 83{,}560 correct $\cdot$ global $z$-score $\mu = -2.41$\,dB, $\sigma = 10.22$\,dB $\cdot$ train 133{,}694 (balanced), test 19{,}662 (natural)};

  \draw[ar] (7.35,-5.60) -- (7.35,-6.25)
        node[midway, right, nt, anchor=west] {one $64 \times 32 \times 1$ patch per phoneme};

  % ===================== BAND 2 : NETWORK =====================
  \draw[band, fill=orange!2] (-0.35,-6.25) rectangle (15.05,-12.35);
  \node[bt] at (-0.15,-6.60)
       {2 --- CNN--BiLSTM CLASSIFIER \normalfont\scriptsize(\S\ref{sec:network-design}, 415{,}810 parameters)};

  \node[nb] (n1) at (1.90,-8.05)  {\textbf{input}\\log-mel patch};
  \node[nb] (n2) at (5.60,-8.05)  {\textbf{conv block 1}\\Conv2D 16 @ $3{\times}3$, ReLU\\BN, MaxPool, Drop 0.3};
  \node[nb] (n3) at (9.30,-8.05)  {\textbf{conv block 2}\\Conv2D 32 @ $3{\times}3$, ReLU\\BN, MaxPool, Drop 0.3};
  \node[nb] (n4) at (13.00,-8.05) {\textbf{reshape}\\TimeDistributed\\Flatten};

  \node[nb] (n5) at (3.70,-10.85) {\textbf{BiLSTM}\\128 units per direction\\Drop 0.3};
  \node[nb] (n6) at (7.90,-10.85) {\textbf{dense}\\64, ReLU\\Drop 0.3};
  \node[nb] (n7) at (12.10,-10.85) {\textbf{output}\\Dense 2, softmax\\$P(\text{error})$};

  \foreach \s/\t in {n1/n2, n2/n3, n3/n4, n5/n6, n6/n7} \draw[ar] (\s) -- (\t);
  % dropped at x = 14.40 so the vertical run clears the shape annotations below
  \draw[ar] (n4.south -| 14.40,0) -- (14.40,-9.75) -- (3.70,-9.75) -- (n5.north);

  \node[nt] at (1.90,-9.20)   {$64{\times}32{\times}1$};
  \node[nt] at (5.60,-9.20)   {$32{\times}16{\times}16$};
  \node[nt] at (9.30,-9.20)   {$16{\times}8{\times}32$};
  \node[nt] at (12.70,-9.20)  {$16{\times}256$};
  \node[nt] at (3.70,-11.90)  {$256$};
  \node[nt] at (7.90,-11.90)  {$64$};
  \node[nt] at (12.10,-11.90) {$2$};

  \draw[ar] (7.35,-12.35) -- (7.35,-13.00);

  % ===================== BAND 3 : TRAINING & EVALUATION =====================
  \draw[band, fill=green!3] (-0.35,-13.00) rectangle (15.05,-15.95);
  \node[bt] at (-0.15,-13.33)
       {3 --- TRAINING \& EVALUATION \normalfont\scriptsize(\S\ref{sec:training-protocol}--\S\ref{sec:metrics})};

  \node[tb] (t1) at (2.55,-14.70) {\textbf{training}\\Adam, categorical cross-entropy,\\batch 32, $\le 30$ epochs,\\early stopping patience 5};
  \node[tb] (t2) at (7.60,-14.70) {\textbf{artefacts}\\trained CNN--BiLSTM\\$+$ normalisation constants $\mu,\sigma$};
  \node[tb] (t3) at (12.65,-14.70) {\textbf{evaluation}\\accuracy, error precision/recall/$F_1$,\\per-phoneme profile,\\control-vs-dysarthric separation};

  \draw[ar] (t1) -- (t2);
  \draw[ar] (t2) -- (t3);

  % ===================== LEGEND =====================
  \fill[red!12, draw=black!65] (0.10,-16.75) rectangle (0.45,-16.51);
  \node[nt, anchor=west] at (0.55,-16.63)
       {contribution of this work: the ground truth is anchored to healthy control speech, not to a pooled percentile};
\end{tikzpicture}}
\caption[Proposed deep-learning pipeline]{The proposed pipeline, end to end.
Band~1 converts a pair of recordings into labelled log-mel patches; band~2 is
the CNN--BiLSTM classifier, annotated with the tensor shape at each stage;
band~3 fits and assesses it. The highlighted stage is the control-baseline
labelling rule that distinguishes this method from the global-percentile
alternative it replaced.}
\label{fig:method-pipeline}
\end{figure}

\section{Dataset}
\label{sec:dataset}

\subsection{Source corpus}

The UASpeech corpus \cite{kim2008dysarthric} was selected for the reason given
in Section~\ref{sec:lit-corpora}: it is the largest publicly available English
corpus that contains \emph{matched} dysarthric and healthy-control productions
of the same isolated words. A noise-reduced distribution of the corpus was used
so that the log-mel deviation between a dysarthric token and its control
counterpart reflects articulatory difference rather than differences in the
recording noise floor.

\subsection{Pairing}

The unit of the experiment is a \emph{pair}: one dysarthric recording and one
control recording of the same target word. A transcript index was prepared in
advance, in which every row supplies four fields: the path to the dysarthric
audio, the path to the control audio, the transcript of the dysarthric
recording and the transcript of the control recording. Ten thousand such pairs
were used.

\subsection{Choice of alignment target}

A pair provides two candidate transcripts, and the choice between them is
consequential.

The transcript obtained from the \emph{dysarthric} recording is a transcription
of what was actually produced --- including its errors. Aligning the dysarthric
audio to that transcript would place phoneme boundaries around the sounds that
were produced, thereby \emph{concealing} the very deviations the system is
supposed to detect. Worse, the two recordings of a pair would then be aligned
to different symbol sequences, and the correspondence between phoneme $i$ of one
and phoneme $i$ of the other --- the correspondence on which the entire
distance-based labelling scheme rests --- would be lost.

The transcript obtained from the \emph{control} recording is a transcription of
clear speech and is therefore the best available proxy for the \emph{intended}
target word. It is used as the alignment target for \emph{both} recordings of
the pair. Phoneme $i$ of the dysarthric alignment and phoneme $i$ of the control
alignment are then, by construction, two attempts at the same intended sound,
and their acoustic distance is a meaningful measure of how far the impaired
attempt fell from a clear realisation. This decision is illustrated in
Figure~\ref{fig:matched-alignment}.

\begin{figure}[H]
\centering
\begin{tikzpicture}[node distance=7mm and 14mm]
  \node[data] (ctrltr) {control transcript \\ \texttt{"shift"}};
  \node[block, below=10mm of ctrltr] (g2p) {espeak-ng G2P};
  \node[data, below=7mm of g2p, fill=yellow!14] (seq)
        {target phoneme sequence \\ \ph{ʃ}\,\ph{ɪ}\,\ph{f}\,\ph{t}};

  \node[data, left=26mm of seq] (dysaud) {dysarthric audio};
  \node[data, right=26mm of seq] (ctrlaud) {control audio};

  \node[block, below=10mm of dysaud] (dysal) {forced align};
  \node[block, below=10mm of ctrlaud] (ctrlal) {forced align};

  \node[data, below=7mm of dysal] (dysseg) {segments $d_1 \dots d_n$};
  \node[data, below=7mm of ctrlal] (ctrlseg) {segments $c_1 \dots c_n$};

  \node[block, below=42mm of seq, fill=red!8] (dist)
        {per-phoneme distance \\ $D_i = \lVert d_i - c_i \rVert_1$};

  \draw[arrow] (ctrltr) -- (g2p);
  \draw[arrow] (g2p) -- (seq);
  \draw[arrow] (seq.west) -- ++(-1.0,0) |- (dysal.east);
  \draw[arrow] (seq.east) -- ++(1.0,0)  |- (ctrlal.west);
  \draw[arrow] (dysaud) -- (dysal);
  \draw[arrow] (ctrlaud) -- (ctrlal);
  \draw[arrow] (dysal) -- (dysseg);
  \draw[arrow] (ctrlal) -- (ctrlseg);
  \draw[arrow] (dysseg.south) |- (dist.west);
  \draw[arrow] (ctrlseg.south) |- (dist.east);
\end{tikzpicture}
\caption[Matched-pair alignment]{Matched-pair alignment. Both recordings are
aligned to the \emph{same} canonical phoneme sequence derived from the control
transcript, guaranteeing that segment $d_i$ and segment $c_i$ are two attempts
at the same intended sound.}
\label{fig:matched-alignment}
\end{figure}

\section{Grapheme-to-Phoneme Conversion}
\label{sec:g2p}

The target word is converted to an IPA phoneme sequence by invoking the
\texttt{espeak-ng} binary directly:

\begin{center}
\texttt{espeak-ng -q --ipa -v en-us <word>}
\end{center}

Three post-processing steps are applied. Primary stress (\texttt{U+02C8}),
secondary stress (\texttt{U+02CC}) and the length mark (\texttt{U+02D0}) are
stripped, because these diacritics are suprasegmental annotations that are not
members of the alignment model's phone vocabulary. The remaining string is then
split into individual symbols, and any symbol not present in the alignment
model's vocabulary is discarded.

The \texttt{espeak-ng} \emph{binary} is invoked through a subprocess rather than
through the \texttt{phonemizer} Python wrapper. This is a practical rather than
a scientific decision: \texttt{phonemizer} copies the shared library into a
temporary directory before loading it, which fails on execution environments
that mount the temporary filesystem \texttt{noexec} --- including the Kaggle
notebook environment used for training. Calling the installed binary in place
avoids the problem entirely, and the identical code path is used in training and
in the deployed service, which eliminates a class of train--serve skew.

A start-up self-check verifies that the toolchain is intact before any data is
processed: the word ``well'' is converted and the result checked to be non-empty
and to intersect the alignment model's vocabulary. The check confirmed
\ph{w}\,\ph{ɛ}\,\ph{l}, all three symbols present in the vocabulary.

\section{Forced Alignment}
\label{sec:alignment}

Alignment uses the multilingual phone recogniser
\pth{facebook/wav2vec2-lv-60-espeak-cv-ft} \cite{xu2022simple} as an emission
model, followed by CTC forced alignment as formalised in
Equation~\eqref{eq:ctc-align}.

The procedure for one recording is as follows.

\begin{enumerate}[leftmargin=*]
  \item \textbf{Load and trim.} The waveform is loaded at 16\,kHz mono, and
        leading and trailing silence is removed with a threshold of
        $20$\,dB below the peak. Trimming is performed \emph{once}, and every
        subsequent operation --- alignment, cropping and, in the deployed
        application, waveform display --- uses the trimmed signal. This
        single-timebase discipline was introduced after an early defect in
        which alignment timestamps computed on the trimmed signal were drawn
        against the untrimmed waveform, displacing the phoneme bands in the
        user interface.
  \item \textbf{Emission.} The waveform is passed through the frozen wav2vec2
        CTC model and a log-softmax is taken over the vocabulary axis, yielding
        $\log P(k \mid t)$ for $T$ frames.
  \item \textbf{Alignment.} \texttt{torchaudio.functional.forced\_align}
        computes the optimal frame path of Equation~\eqref{eq:ctc-align} for the
        target token sequence, and \texttt{merge\_tokens} collapses that path
        into one span per target symbol, each with a mean posterior confidence.
  \item \textbf{Timebase conversion.} Frame indices are converted to seconds
        with the ratio
        $r = N_{\text{samples}} / (T \cdot f_s)$,
        giving $(p_i, t^{\text{start}}_i, t^{\text{end}}_i, s_i)$ for every
        phoneme, where $s_i \in [0,1]$ is the alignment confidence.
\end{enumerate}

Three pairs out of ten thousand failed at this stage, all with the same
underlying cause: the target sequence was longer than the number of emission
frames, which makes the constrained path set $\mathcal{B}^{-1}(\mathbf{y})$
empty. Such pairs were discarded.

\section{Feature Extraction}
\label{sec:features}

Each aligned phoneme is converted to a fixed-size log-mel patch.

\textbf{Context padding.} The segment is widened by $20$\,ms on each side. Stop
consonants in particular are only interpretable together with their closure and
release, and the alignment boundary tends to fall inside the acoustic event
rather than around it. The padded interval is clipped to the bounds of the
waveform.

\textbf{Log-mel spectrogram.} A mel spectrogram with $M = 64$ mel bands is
computed over the cropped clip with the default \texttt{librosa} analysis
parameters (2048-sample window, 512-sample hop) and converted to decibels
according to Equation~\eqref{eq:logmel}.

\textbf{Fixed width.} The classifier requires a constant input shape. Patches
wider than $F = 32$ frames are truncated to the first 32; narrower patches are
zero-padded on the right. Truncation is preferred to resampling because
duration is itself disrupted in dysarthria
(Section~\ref{sec:lit-dysarthria}) and rescaling the time axis would remove
that information; the 32-frame window at the stated hop covers approximately
$1.0$\,s of audio, which exceeds any plausible single-phoneme duration, so
truncation is in practice rare.

\textbf{Normalisation.} A single global $z$-score is applied across the entire
feature tensor,
\begin{equation}
\tilde{X} = \frac{X - \mu}{\sigma + \epsilon},
\qquad
\mu = \operatorname{mean}(X),\quad \sigma = \operatorname{std}(X),
\label{eq:zscore}
\end{equation}
with $\mu$ and $\sigma$ computed once over the training corpus and
\emph{persisted}. The measured values were $\mu = -2.41$\,dB and
$\sigma = 10.22$\,dB. Because the log-mel representation is in decibels and
therefore level-dependent, these two scalars are as much a part of the trained
model as the network weights; the deployed service reproduces them exactly
(Section~\ref{sec:impl-normalisation}).

The complete parameter set is given in Table~\ref{tab:feature-params}. Every one
of these values is duplicated in the deployed service's configuration module and
must not diverge.

\begin{table}[H]
\centering
\caption[Audio and feature extraction parameters]{Audio and feature extraction
parameters, shared by training and inference.}
\label{tab:feature-params}
\begin{tabular}{@{}llp{6.6cm}@{}}
\toprule
\textbf{Parameter} & \textbf{Value} & \textbf{Rationale} \\
\midrule
Sample rate            & 16\,000\,Hz & Native rate of the wav2vec2 aligner \\
Channels               & mono        & Single articulatory source \\
Silence trim threshold & 20\,dB      & Removes leading/trailing silence \\
Mel bands $M$          & 64          & Frequency resolution vs.\ patch size \\
Fixed frames $F$       & 32          & Constant classifier input width \\
Context padding        & 0.02\,s     & Preserves closure/release of stops \\
Patch shape            & $64\times32\times1$ & Single-channel image \\
Global mel mean $\mu$  & $-2.41$\,dB & Persisted normalisation constant \\
Global mel std $\sigma$& $10.22$\,dB & Persisted normalisation constant \\
\bottomrule
\end{tabular}
\end{table}

\section{Label Derivation}
\label{sec:labelling}

No clinician-annotated per-phoneme labels exist for this corpus
(Section~\ref{sec:lit-corpora}), so labels must be derived from the paired
structure of the data. Two strategies were implemented. The first was adopted
initially and abandoned; the second is the method of this dissertation.

\subsection{Phoneme deviation measure}

Both strategies rest on the same underlying quantity. For the $i$-th phoneme of
a pair, let $\mathbf{d}_i$ and $\mathbf{c}_i$ be the log-mel patches of the
dysarthric and matched control segments. The deviation is the mean absolute
difference in decibels over the overlapping region,
\begin{equation}
D_i = \frac{1}{M \cdot F'} \sum_{m=1}^{M} \sum_{t=1}^{F'}
      \bigl| \mathbf{d}_i(m,t) - \mathbf{c}_i(m,t) \bigr|,
\qquad F' = \min\bigl(F_{d_i}, F_{c_i}\bigr).
\label{eq:distance}
\end{equation}
The $L_1$ norm is used rather than the $L_2$ norm because it is less dominated
by isolated large-magnitude bins, which in log-mel space are frequently caused
by low-energy frames rather than by articulatory difference. The measure is
computed on the \emph{unnormalised} decibel patches so that it carries a
physical interpretation: $D_i$ is an average difference in decibels.

Over the corpus, $D$ ranged from $0.13$\,dB to $6.83$\,dB with a median of
$0.76$\,dB.

\subsection{Method 1: global percentile labelling}
\label{sec:method1}

The first strategy pooled every phoneme instance --- dysarthric and control
alike --- into a single distribution of $D$, and declared the top
$\rho = 30\%$ of that pooled distribution to be errors:
\begin{equation}
y_i^{(1)} = \mathbb{1}\bigl[\,D_i \geq \operatorname{Perc}_{100-\rho}\bigl(\{D_j\}_{j \in \text{all}}\bigr)\,\bigr].
\label{eq:method1}
\end{equation}

The appeal of Equation~\eqref{eq:method1} is that it requires no assumption
about which speakers are impaired, and it produces a class balance of exactly
$\rho$ by construction, which is convenient for training.

That convenience is the defect. Because the percentile is computed over the
pooled distribution, the criterion is not ``deviates from healthy speech'' but
``is among the most deviant $30\%$ of everything''. Healthy control phonemes
occupy the lower part of the distribution but not exclusively so, and a
non-trivial share of them therefore exceeds the pooled cutoff and is labelled
erroneous \emph{even though it was produced by a healthy speaker with no
articulation impairment}. The label noise this introduces is systematic and is
concentrated precisely on the examples that ought to define the correct class.

The consequence was observed directly at inference time. A model trained under
Equation~\eqref{eq:method1} learned the base rate implied by the labelling
scheme and flagged roughly that proportion of phonemes in \emph{any} recording.
Presented with a clear control recording, it reported errors throughout. The
model was therefore unable to perform the one discrimination the system exists
to perform --- separating clear from impaired speech --- and the strategy was
abandoned.

\subsection{Method 2: control-baseline labelling}
\label{sec:method2}

Method~2 changes the definition of correctness from a statistical one to a
referential one: \emph{correct} means \emph{like a healthy speaker producing
this sound}.

Let $\mathcal{C}$ denote the set of control phoneme instances and $\mathcal{D}$
the set of dysarthric phoneme instances. The labelling rule is
\begin{equation}
y_i^{(2)} =
\begin{cases}
0, & i \in \mathcal{C},\\[4pt]
\mathbb{1}\bigl[\,D_i \geq \tau\,\bigr], & i \in \mathcal{D},
\end{cases}
\qquad
\tau = \operatorname{Perc}_{100-\rho}\bigl(\{D_j\}_{j \in \mathcal{D}}\bigr).
\label{eq:method2}
\end{equation}

The first branch is unconditional: a control instance is correct by definition.

Three properties follow directly from Equation~\eqref{eq:method2} and together
constitute the contribution.

\begin{enumerate}[leftmargin=*]
  \item \textbf{Control phonemes cannot be labelled as errors.} The first branch
        is unconditional. There is no value of $D$, no percentile and no
        threshold that can move a control instance into the error class. The
        systematic mislabelling of Method~1 is eliminated by construction rather
        than mitigated by tuning.
  \item \textbf{The threshold is estimated only over impaired data.} $\tau$ is
        the $(100-\rho)$-th percentile of $\{D_j\}_{j \in \mathcal{D}}$, not of
        the pooled set. The cutoff therefore answers the question ``which
        dysarthric productions deviate most from their own healthy
        counterparts'', which is the clinically meaningful question.
  \item \textbf{Well-produced dysarthric phonemes are retained as correct.} The
        second branch is a comparison, not a quota over the whole corpus. A
        dysarthric speaker who produces a sound cleanly yields $D_i < \tau$ and
        that instance enters training as a positive example of correctness.
        This matters both for fairness --- the model does not learn ``impaired
        speaker $\Rightarrow$ error'' --- and for clinical utility, since
        recognising a successful production is what makes progress visible.
\end{enumerate}

Additionally, every control instance is \emph{retained in the training set} as
an explicit correct example rather than discarded. The control instances are
the anchors that define the correct class; removing them would leave the model
with only the low-deviation dysarthric instances as its notion of correctness,
which is a weaker and noisier definition.

The value $\rho = 30\%$ was chosen as a compromise: a smaller $\rho$ yields
cleaner error labels but too few positive examples of the minority class, while
a larger $\rho$ dilutes the error class with mildly deviant productions. With
$\rho = 30\%$ the resulting cutoff was $\tau = 1.20$\,dB.

Figure~\ref{fig:labelling-comparison} contrasts the two schemes, and
Table~\ref{tab:labelling-comparison} summarises the differences.

\begin{figure}[H]
\centering
\begin{tikzpicture}
\begin{groupplot}[
  group style={group size=2 by 1, horizontal sep=16mm},
  width=0.48\textwidth, height=5.2cm,
  xlabel={Deviation $D$ (dB)},
  xmin=0, xmax=4, ymin=0, ymax=1.25,
  ytick=\empty, xtick={0,1,2,3,4},
  tick label style={font=\scriptsize}, label style={font=\scriptsize},
  title style={font=\scriptsize},
  legend style={font=\tiny, at={(0.97,0.97)}, anchor=north east, draw=none,
                fill=none},
]

\nextgroupplot[ylabel={density}, title={Method 1: pooled percentile}]
\addplot[smooth, thick, draw=green!55!black, fill=green!25, fill opacity=0.45,
         domain=0:4, samples=120]
  {exp(-((x-0.55)^2)/(2*0.30^2))};
\addlegendentry{control}
\addplot[smooth, thick, draw=red!65!black, fill=red!22, fill opacity=0.45,
         domain=0:4, samples=120]
  {0.85*exp(-((x-1.30)^2)/(2*0.70^2))};
\addlegendentry{dysarthric}
\draw[very thick, black, dashed] (axis cs:1.05,0) -- (axis cs:1.05,1.15);
\node[font=\tiny, anchor=south] at (axis cs:1.05,1.16) {pooled cutoff};
\node[font=\tiny, anchor=west, text=red!60!black, align=center]
     at (axis cs:1.15,0.55) {control tail\\mislabelled};

\nextgroupplot[ylabel={}, title={Method 2: control-baseline}]
\addplot[smooth, thick, draw=green!55!black, fill=green!25, fill opacity=0.45,
         domain=0:4, samples=120]
  {exp(-((x-0.55)^2)/(2*0.30^2))};
\addlegendentry{control (exempt)}
\addplot[smooth, thick, draw=red!65!black, fill=red!22, fill opacity=0.45,
         domain=0:4, samples=120]
  {0.85*exp(-((x-1.30)^2)/(2*0.70^2))};
\addlegendentry{dysarthric}
\draw[very thick, black, dashed] (axis cs:1.20,0) -- (axis cs:1.20,1.15);
\node[font=\tiny, anchor=south, align=center] at (axis cs:1.20,1.16)
     {$\tau$ from dysarthric only};
\end{groupplot}
\end{tikzpicture}
\caption[Method 1 versus Method 2 labelling]{Schematic comparison of the two
labelling strategies. Under Method~1 the cutoff is estimated over the pooled
distribution, so the upper tail of the control distribution is forced into the
error class. Under Method~2 the cutoff $\tau$ is estimated over the dysarthric
distribution alone and control instances are exempt from thresholding
altogether.}
\label{fig:labelling-comparison}
\end{figure}

\begin{table}[H]
\centering
\caption[Comparison of labelling strategies]{Comparison of the two labelling
strategies.}
\label{tab:labelling-comparison}
\small
\begin{tabularx}{\textwidth}{@{}p{4.1cm}XX@{}}
\toprule
\textbf{Aspect} & \textbf{Method 1 (global percentile)} & \textbf{Method 2 (control baseline)} \\
\midrule
Definition of ``correct'' & Below a pooled statistical cutoff & Produced by a healthy control speaker \\
Threshold estimated over  & All phonemes (control + dysarthric) & Dysarthric phonemes only \\
Control instances         & Thresholded like any other; upper tail becomes ``error'' & Labelled correct unconditionally; retained as anchors \\
Well-produced dysarthric  & May be forced into the error class to satisfy the quota & Correctly labelled as correct \\
Class balance             & Exactly $\rho$ by construction & Emergent (15.0\% error observed) \\
Failure mode              & Flags errors in clear speech at inference & --- \\
\bottomrule
\end{tabularx}
\end{table}

\subsection{Resulting dataset composition}

Applying Equation~\eqref{eq:method2} to the aligned corpus yielded
$49{,}153$ control instances and $49{,}153$ dysarthric instances, a total of
$98{,}306$ labelled phoneme patches. Of the dysarthric instances, $14{,}746$
exceeded $\tau$ and were labelled as errors and $34{,}407$ were labelled
correct. Combined with the control instances, the corpus contains $14{,}746$
error and $83{,}560$ correct instances --- a 15.0\%/85.0\% split. This imbalance
is an honest reflection of the phenomenon (most attempted phonemes, even by
impaired speakers, are not gross errors) and is handled at the training stage
rather than by distorting the labels.

\section{Data Splitting}
\label{sec:split}

The corpus was split 80\%/20\% into training and held-out sets with
stratification on the label, under a fixed random seed of 42.

\textbf{The split precedes class balancing.} An earlier iteration of the
pipeline balanced the classes by upsampling and \emph{then} split, with the
result that duplicated copies of the same minority instance appeared on both
sides of the split. The held-out metrics obtained under that ordering were
inflated and meaningless. In the final pipeline, the split is performed first
and balancing is applied to the training partition only; the held-out partition
retains the natural 15\%/85\% distribution.

\textbf{Balancing.} Within the training partition the minority class is
upsampled with replacement to the size of the majority class, and the result is
shuffled. This yields a training set of $133{,}694$ instances at exactly 50/50
balance, and a held-out set of $19{,}662$ instances at the natural distribution.

\textbf{A stated limitation.} The held-out partition was also supplied to Keras
as the validation set for early stopping, so the reported held-out metrics are
obtained on data that participated in the stopping decision. This is a genuine
methodological weakness; it is quantified and discussed in
Section~\ref{sec:limitations} rather than concealed. A three-way split with a
separate validation partition is the correct design and is recommended as
future work.

\section{Model Training Protocol}
\label{sec:training-protocol}

All candidate architectures were trained under the identical protocol given in
Table~\ref{tab:hyperparams}, so that the architecture comparison of
Section~\ref{sec:arch-design} is controlled.

\begin{table}[H]
\centering
\caption[Training hyperparameters]{Training hyperparameters, held constant
across all candidate architectures.}
\label{tab:hyperparams}
\small
\begin{tabular}{@{}llp{5.2cm}@{}}
\toprule
\textbf{Hyperparameter} & \textbf{Value} & \textbf{Note} \\
\midrule
Loss                 & Categorical cross-entropy & Two-unit softmax output \\
Optimiser            & Adam                      & $\eta = 10^{-3}$, $\beta_1 = 0.9$, $\beta_2 = 0.999$ \\
Batch size           & 32                        & \\
Maximum epochs       & 30                        & \\
Early stopping       & patience 5                & On \texttt{val\_loss}; best weights restored \\
Dropout              & 0.3                       & After each block and the dense layer \\
Random seed          & 42                        & Split, upsampling and shuffle \\
Class balancing      & Upsampling                & Training partition only \\
\bottomrule
\end{tabular}
\end{table}

Dropout is applied at a relatively high rate of $0.3$ throughout because the
labels are proxy labels derived from an acoustic distance measure and therefore
carry appreciable noise; the regularisation is intended to discourage the
network from fitting that noise. Recurrent dropout is deliberately \emph{not}
used, because enabling it disqualifies the layer from the fused cuDNN kernel and
increases training time by roughly an order of magnitude for no measured
benefit.

\section{Architecture Comparison Design}
\label{sec:arch-design}

Five architectures were compared, all sharing the same convolutional front end
and differing only in the temporal-modelling stage:

\begin{enumerate}[leftmargin=*]
  \item \textbf{CNN} --- convolutional blocks followed directly by the dense
        classifier (no recurrent stage). This is the ablation that quantifies
        the contribution of temporal modelling.
  \item \textbf{CNN--LSTM} --- unidirectional LSTM.
  \item \textbf{CNN--GRU} --- unidirectional GRU.
  \item \textbf{CNN--BiLSTM} --- bidirectional LSTM.
  \item \textbf{CNN--BiGRU} --- bidirectional GRU.
\end{enumerate}

\textbf{Selection criterion.} Accuracy is an inadequate criterion for this
task. Under the natural class distribution a model that predicts ``correct''
unconditionally achieves 85\% accuracy while detecting nothing. The criterion
adopted is \textbf{recall on the error class}, with error-class $F_1$ as the
tie-breaker, on the grounds that the two error types carry asymmetric costs in
a screening application. A false negative --- a genuinely unclear sound reported
as fine --- silently withholds the feedback the tool exists to provide. A false
positive --- a correct sound flagged for attention --- costs the user an
unnecessary repetition, which is the ordinary business of therapy practice. The
system is explicitly positioned as a screening instrument, and screening
instruments are tuned for sensitivity. The user-adjustable threshold described
in Section~\ref{sec:design-threshold} allows this operating point to be moved
by the user if a more conservative behaviour is preferred.

\section{Evaluation Metrics}
\label{sec:metrics}

With the error class as the positive class, the standard quantities are used:
\begin{align}
\text{Accuracy}  &= \frac{TP + TN}{TP + TN + FP + FN}, &
\text{Precision} &= \frac{TP}{TP + FP}, \label{eq:acc-prec}\\[4pt]
\text{Recall}    &= \frac{TP}{TP + FN}, &
F_1              &= \frac{2 \cdot \text{Precision} \cdot \text{Recall}}
                          {\text{Precision} + \text{Recall}}. \label{eq:rec-f1}
\end{align}

In addition to these, two task-specific evaluations are reported.

\textbf{Control-versus-dysarthric separation.} The mean fraction of phonemes
flagged per recording is computed separately over control recordings and over
dysarthric recordings, and the gap between them is reported. This measures
directly the property that Method~1 failed to deliver, and it is the evaluation
that the labelling redesign was aimed at.

\textbf{Per-phoneme error profile.} The predicted error rate is computed for
each IPA symbol, restricted to symbols occurring at least 100 times so that the
estimate is not dominated by sampling noise. The resulting ranking is compared
against articulatory patterns reported in the clinical literature, as an
external plausibility check on what the model has learned.

\section{Tools and Experimental Environment}
\label{sec:tools}

Training was performed in a Kaggle notebook on an NVIDIA Tesla T4 GPU with
CUDA-accelerated TensorFlow. The deployed inference service targets CPU-only
commodity hardware. The principal software components are listed in
Table~\ref{tab:tools}.

\begin{table}[H]
\centering
\caption[Tools and libraries]{Principal tools and libraries.}
\label{tab:tools}
\begin{tabular}{@{}llp{6.0cm}@{}}
\toprule
\textbf{Component} & \textbf{Version} & \textbf{Role} \\
\midrule
Python              & 3.11            & Implementation language \\
TensorFlow / Keras  & 2.16 / 3        & Classifier training and inference \\
PyTorch             & 2.2.2           & Backend for the aligner \\
torchaudio          & 2.2.2           & CTC forced alignment \\
transformers        & 4.44.2          & wav2vec2 and Whisper model loading \\
librosa             & 0.10.2          & Audio I/O, trimming, mel spectrograms \\
NumPy / SciPy       & 1.26 / 1.13     & Numerical computation \\
scikit-learn        & ---             & Splitting, resampling, metrics \\
espeak-ng           & 1.50            & Grapheme-to-phoneme conversion \\
FastAPI / Uvicorn   & 0.115 / 0.34    & Inference service \\
React / Vite        & 18.3 / 6.0      & Web client \\
Docker              & ---             & Packaging and deployment \\
\bottomrule
\end{tabular}
\end{table}

\section{Ethical Considerations}
\label{sec:ethics}

\textbf{Data.} All model development used the publicly released UASpeech
corpus, whose contributors gave informed consent under the original collection
protocol and whose recordings are distributed in de-identified form. No new
recordings of human participants were collected for this dissertation, and no
clinical trial was conducted.

\textbf{Privacy in deployment.} The deployed service processes uploaded audio in
memory and does not persist it. There are no user accounts, no session history
and no analytics on recorded content. This is a deliberate design choice for a
system whose input is a biometric signal from a potentially vulnerable
population.

\textbf{Framing of the output.} The system's output could be misread as a
clinical assessment. Three mitigations are applied. The interface presents a
clarity score and highlights sounds for practice, and never presents a
diagnosis or a severity label. The documentation states plainly that the system
is a screening signal and not a clinical instrument. And the sensitivity
threshold is exposed to the user, which makes visible the fact that the verdict
is a tunable decision rather than an objective fact about the speech.

\textbf{Bias.} The model was trained on the speakers present in UASpeech, whose
dysarthria arises predominantly from cerebral palsy, and on American English.
Performance on other aetiologies, other accents, other age groups and other
languages is unmeasured and should not be assumed. This is stated as a
limitation in Section~\ref{sec:limitations} and as a priority for future work.
